This chapter presented a Root Cause Diagnosis (RCD) approach for identifying network issues in Cloud VR sessions over Wi-Fi networks, utilizing time series KPIs collected from access points. By employing a two-stage framework, we demonstrated the effectiveness of our approach compared to traditional time series classification methods. Our proposed architecture, which integrates contrastive learning into the anomaly detection process, has shown significant improvements in both identifying anomalies and diagnosing the root causes of Cloud VR performance issues. This provides a good foundation for future research in real-time diagnostics for cloud-based VR applications. One key strength of this approach is its balance between training time and inference speed, making it ideal for real-time diagnostics in dynamic network environments. Moreover, the two-stage design enhances efficiency by isolating retraining to the supervised classifier, thus avoiding full model retraining when new causes are introduced.

Despite the promising results, there are several areas that warrant further exploration and improvement. First, while our model performs well in detecting clear anomaly patterns such as interference, its sensitivity to more nuanced variations—particularly in attenuation scenarios—needs to be enhanced. Future work will focus on improving the model’s ability to detect subtle transitions between normal and degraded states. Second, scaling the solution to larger and more complex datasets is a priority. Our current framework has been tested on a controlled Cloud VR testbed with only two types of impairments. To improve its robustness, future research should include additional Wi-Fi impairments such as network congestion, hidden terminal issues, and non-Wi-Fi interference. Expanding the test environment to include more diverse real-world conditions, such as home networks where multiple sources of impairments may coexist, will also offer valuable insights. Finally, extending this two-stage model to a multi-modal diagnostic approach could provide a more comprehensive view of the root cause diagnosis process. By incorporating additional data sources, such as application-level performance metrics or raw network PCAP data, the system could offer even more accurate and proactive detection of network impairments. These enhancements will be crucial for addressing the growing demands of next-generation low-latency applications like Cloud VR.